\section{Introduction}
\label{sec:introduction}

The spacetime Penrose inequality compares the invariant ADM mass of
initial data with the least area needed to enclose a trapped boundary.
Its model equality cases are spacelike slices of the
Schwarzschild--Tangherlini spacetime \cite{Tangherlini1963}. Such a slice
need not be time-symmetric. Admissible asymptotically boosted slices can
have different second fundamental forms and ADM energy--momentum vectors
while retaining the same invariant mass. The rigidity question is therefore
whether equality determines the original metric and second fundamental
form as the induced data of a Schwarzschild--Tangherlini slice.

We prove this conclusion in every spatial dimension $n\ge3$ under
strong differentiated decay, and under weaker decay in dimensions
three and four. The embedding includes the original horizon in regular
spacetime coordinates. In dimension three the finite-component theorem
also proves that equality forces the horizon to be connected. The
hypotheses and converse statements differ between the three classes;
we describe those distinctions before explaining the proof.

\subsection{The equality statements and their scope}

Write $(E,P)$ for the original ADM energy--momentum vector and
$m=(E^2-|P|^2)^{1/2}$ for its invariant mass. Let $a(g)$ be the
infimum of the areas of all full enclosing cuts, including every
component and every portion coincident with the original boundary $S$.
The numerical inequality is
\[
 m\ge F_n(a(g)),\qquad
 F_n(a)=\frac12\left(\frac{a}{\omega_{n-1}}\right)^{(n-2)/(n-1)},
 \qquad \omega_{n-1}=|\mathbb S^{n-1}|.
\]
All our equality classes have smooth orientable data satisfying the
dominant energy condition, integrable constraint densities and finite
ADM limits. Their boundary is marginally outer trapped and outer
area-minimizing, so that $a(g)=A:=|S|_g>0$. The chosen closed exterior
is complete with its boundary included and has compact complement to
one asymptotically flat end. No spin, vacuum, maximality or
stationary-development assumption is imposed.

Table~\ref{tab:intro:branches} records the remaining distinctions.
Here $O_j(r^{-a})$ includes the coordinate derivative estimates through
order $j$; the underlying data are smooth. Wholly-interior outermostness
excludes any smooth weakly outer trapped enclosing cut entirely in the
open exterior, even if that cut is disconnected. Partial-coincidence
outermostness also excludes cuts retaining some entire original
boundary components and replacing others by interior components. The
precise definitions are given with the theorems.

\begin{table}[htbp]
\centering
\small
\renewcommand{\arraystretch}{1.13}
\begin{tabular}{@{}p{.14\textwidth}p{.28\textwidth}p{.50\textwidth}@{}}
\toprule
Class & Differentiated asymptotics & Horizon, ambient hypothesis and conclusion \\
\midrule
$n\ge3$ &
$g-\delta=O_6(r^{-q})$,\newline
$K=O_5(r^{-1-q})$,\newline
$\frac{n-2}{2}<q<n-2$ &
Connected $S$; wholly-interior outermostness; no ambient extension
required. Equality and the converse hold in this class
(Theorems~\ref{thm:rigidity} and~\ref{thm:converse}). \\
\addlinespace
$n=3$ &
$g-\delta=O_2(r^{-q})$,\newline
$K=O_1(r^{-1-q})$,\newline
$q>\frac12$ on each\newline ambient end &
The exterior lies in complete boundaryless data with finitely many AF
ends. For finitely many components, partial-coincidence outermostness
gives rigidity and connectedness; for connected $S$, wholly-interior outermostness
suffices. These are forward statements
(Theorems~\ref{q3:intro:rigidity} and~\ref{q3:intro:connected-rigidity}). \\
\addlinespace
$n=4$ &
$g-\delta=O_2(r^{-q})$,\newline
$K=O_1(r^{-1-q})$,\newline
$q>1$ &
Connected $S$; wholly-interior outermostness; no ambient extension
required. Equality and the converse hold in this weak-decay class
(Theorem~\ref{q4:thm:rigidity}). \\
\bottomrule
\end{tabular}
\caption{The three equality classes. Common energy, area and completeness
hypotheses are stated above; the cited theorems give the full assumptions.}
\label{tab:intro:branches}
\end{table}

Future timelikeness $E>|P|$ is assumed in the strong and
three-dimensional statements. In dimension four it follows from the
numerical theorem and positive enclosing area. In the
three-dimensional theorem the ambient completeness assumption ensures
completeness of the closed chosen exterior; the conclusion concerns
only that exterior. No theorem identifies a region behind $S$.

In each branch, equality embeds the entire original exterior, with
its given $g$ and $K$, into the regular extension of the exterior of
mass $m$. The boundary is a full smooth section of the future horizon
or the bifurcation sphere, and the chosen end approaches spatial
infinity. The converses include asymptotically boosted slices and
compute the invariant mass from the original ADM integrals. The
strong class contains examples with $K\not\equiv0$. In the
three-dimensional finite-component class,
Corollary~\ref{q3:rigidity:finite-strictness} gives the further consequence
\[
 S\text{ disconnected}\quad\Longrightarrow\quad
 \sqrt{E^2-|P|^2}>\sqrt{A/(16\pi)}.
\]
This is qualitative strictness, with no asserted uniform positive
lower bound for the deficit.

Two numerical inputs have distinct roles. The spacetime theorem of
\emph{The spacetime Penrose inequality and enclosing area}
\cite[Theorems~1.3, 7.2 and 8.2]{CompanionNumerical} applies to the
nearby weakly trapped data used in the variations; they need not satisfy
the original pair's outermostness or area-minimizing hypotheses. The
Riemannian theorem of \emph{Conformal flow and the Riemannian Penrose
inequality with minimizing frontiers}
\cite[Theorem~1.1]{CompanionRiemannian} enters later in the
all-dimensional classification. We use its numerical bound to prove
nonnegative mass and then rigidity for a rough conformal double.
Neither input supplies the original-data equality argument.

\subsection{Why equality must be varied on the original data}

The numerical spacetime proof constructs auxiliary metrics and passes
to inequalities in a prescribed order. Equality in its final bound
does not identify the original pair $(g,K)$ with any one auxiliary
metric. We instead use the inequality on admissible variations of
$(g,K)$ itself: equality makes the mass--area deficit
$m(g,K)-F_n(a(g))$ a constrained minimum.

There are two obstacles to extracting stationary equations from this
minimum. A least-area enclosing surface can change when the metric
changes. Even when $S$ minimizes area, its own first variation need
not be the derivative of the infimum. Moreover, separation of the
first variations produces measures on all minimizing enclosures.
Interior sheets can therefore contribute an area source to the
adjoint equations. This source must be removed before the equations
of a stationary spacetime become available.

The geometric step attaches an inner side to $S$ and minimizes full
perimeter among sets containing it. Mean-convex spheres confine
minimizers for asymptotically flat metric paths. The broader paths
needed in dimension three use uniform metric comparison and density
bounds instead. Perimeter convergence and continuity for vector
measures control the tangent-plane integrals and give the right
metric derivative
\[
 a'(g;h)=\min_C\frac12\int_{\partial^*C}
                    \operatorname{tr}_{T\partial^*C}h\,dA_g,
\]
where $C$ ranges over all original minimizing enclosures. Thus no
differentiable choice of a minimizing surface is needed.

Positive separation then produces a lapse $u$, a shift $X$ and a
positive area measure. Causality takes the form $u\ge|X|$, while the
Hessian equation initially has a measure source normal to the
minimizing sheets. Compact normal tests show that a typical free
sheet has both zero mean curvature and zero tangential trace of $K$.
In dimensions at most seven it is already smooth, so outermostness
removes it. In higher dimensions, atomic sheets are totally geodesic;
neighboring nonatomic sheets produce a common positive Jacobi field.
Its behavior near singular points, together with a BV zero-set
argument, gives a curvature bound that removes the remaining
singularities. Smooth outermostness can then be applied.

After localization the interior adjoint is homogeneous. Its stationary
development still admits a null-dust alternative. The common
stationary argument, Proposition~\ref{prop:static:common-base}, removes
the twist and constructs the static base without presupposing a
regular interior null set. It retains possible complete ends arising
there. The three classifications subsequently resolve these ends in
different ways, as shown in Figure~\ref{fig:intro:proof}. In particular,
the weak-decay branches derive the harmonic-coordinate estimates
needed for compactification; they do not assume strong differentiated
decay for the original data.

\begin{figure}[htbp]
\centering
\begin{tikzpicture}[
 >=Latex,
 box/.style={draw,rounded corners=2pt,align=center,inner sep=5pt,
             font=\small},
 branch/.style={box,text width=.265\textwidth,minimum height=24mm},
 node distance=6mm]
\node[box,text width=.84\textwidth] (variation)
 {Spacetime numerical inequality on variations of the original data\\
  $\longrightarrow$ full-area envelope and causal multipliers};
\node[box,text width=.84\textwidth,below=of variation] (localization)
 {Normal tests and outermostness localize the area source to $S$};
\node[box,text width=.84\textwidth,below=of localization] (base)
 {Homogeneous adjoint $\longrightarrow$ stationary field $\longrightarrow$
  static base\\Possible interior null ends are retained};
\node[branch,below=9mm of base,xshift=-.33\textwidth] (strong)
 {$n\ge3$, strong decay\\Circle completion\\
  Rough conformal double\\Riemannian mass bound};
\node[branch,below=9mm of base] (three)
 {$n=3$, weak decay\\Complete conformal double\\
  Spinor rigidity with other ends};
\node[branch,below=9mm of base,xshift=.33\textwidth] (four)
 {$n=4$, weak decay\\$C^2$ compactification\\
  Nonspin nonnegative mass};
\draw[->] (variation)--(localization);
\draw[->] (localization)--(base);
\draw[->] (base.south)--(three.north);
\draw[->] (base.south) -- ++(0,-4mm) -| (strong.north);
\draw[->] (base.south) -- ++(0,-4mm) -| (four.north);
\end{tikzpicture}
\caption{The common variational and stationary steps precede three
complete classification arguments. Each branch proves the rigidity
needed for its conformal double and then recovers the original slice
through the horizon. Numerical mass bounds are used only in the
indicated roles.}
\label{fig:intro:proof}
\end{figure}

\subsection{Mathematical context and attribution}

Penrose's mass--area question arose from black-hole collapse and cosmic
censorship \cite{Penrose1973}. The area in an initial-data formulation
requires care. Ben-Dov constructed spherically symmetric data satisfying
the dominant energy condition for which the bound using the apparent
horizon's own area fails; he distinguished that area from the minimum
area needed to enclose it \cite[Section~1]{BenDov2004}. This distinction
explains the use of $a(g)$ in the numerical inequality and the separate
outer area-minimizing assumption in our equality classes.

For time-symmetric data, Huisken and Ilmanen's weak inverse mean
curvature flow proves the three-dimensional bound associated with a
connected horizon, while Bray's conformal flow gives the sharp bound
for the total area of a possibly disconnected horizon
\cite{HuiskenIlmanen2001,Bray2001}. Bray and Lee develop the conformal
flow in dimensions below eight \cite{BrayLee2009}. Minimizing hulls
and smooth-obstacle regularity are central to the inverse mean
curvature approach. Here perimeter compactness is combined with
Reshetnyak's continuity principle to differentiate the enclosing
infimum as the metric varies \cite{Reshetnyak1968,Spector2011}.

Beyond time symmetry, Bray and Khuri's generalized Jang construction
is an important equality precedent. In their spherically symmetric
setting they recover the original metric and second fundamental form
as a Schwarzschild spacelike embedding
\cite[Section~4]{BrayKhuri2010}. Their proposed nonspherical coupled
system is conditional on its existence and regularity
\cite[Section~5]{BrayKhuri2010}. The original-slice conclusion, rather
than vanishing of $K$, is thus already integral to the earlier
spacetime formulation. The argument here obtains that conclusion by
varying the original mass--area deficit after the numerical theorem
has been established.

Bartnik's variational program relates constrained gravitational
energy to stationary geometry. In his boundaryless asymptotically
flat three-dimensional phase space, critical points of the ADM energy
on a fixed constraint level set are characterized by the adjoint constraint
equations; future-timelike ADM vectors give the corresponding
invariant-mass formulation
\cite[Theorem~6.1 and Corollary~6.2]{Bartnik2005PhaseSpace}.
The Killing-initial-data framework and asymptotic Killing-field
estimates of Beig and Chru\'sciel explain how lapse--shift equations
encode a spacetime symmetry
\cite{BeigChrusciel1996,BeigChrusciel1997KID}.
For the dominant energy condition, Corvino and Huang introduce a
modified constraint operator that accounts for the metric dependence
of the momentum norm \cite[Section~2.2]{CorvinoHuang2020}.
Huang and Lee relate its adjoint, together with a null-vector identity,
to Killing developments with null-fluid stress
\cite[Section~6.1]{HuangLee2024Bartnik}.

Hirsch and Huang obtain a strongly stationary vacuum development of
the interior for local ADM-mass minimizers with compact boundary,
fixed Bartnik data, the dominant energy condition and $E>|P|$, under
their stated decay and regularity hypotheses
\cite[Theorem~7]{HirschHuang2025}. Fixing those boundary data fixes
$|S|$, but need not fix the least area of all enclosing cuts. Our
constrained minimum is therefore a different variational problem:
it retains $a(g)$ and the interior source generated by its minimizing
enclosures. Their strong maximum principle rules out an interior null
point for a causal Killing field $Y$ that is timelike near infinity
and satisfies $\operatorname{Ric}(Y,Y)\le0$
\cite[Theorem~4]{HirschHuang2025}. Here localization first produces the
homogeneous adjoint; the stationary construction retains its possible
null stress, and the twist argument treats the null set directly.

Static classification has its own ancestry. In three spatial
dimensions, Bunting and Masood-ul-Alam used positive mass to remove
the connected-horizon assumption in the asymptotically Euclidean
static vacuum problem \cite{BuntingMasoodUlAlam1987}.
Gibbons, Ida and Shiromizu developed the conformal construction in
higher dimensions \cite[Section~2]{GibbonsIdaShiromizu2002}.
Hwang's related rigidity theorem concerns complete spin Ricci-flat
manifolds with prescribed asymptotically flat behavior and an
appropriate circle action with fixed points
\cite{Hwang1998}; this is also a precedent for the circle-extension
setting used here. For the present data, we establish the completeness
of each double and the required horizon and compactified-point regularity.

The mass inputs are accordingly different in the three branches.
Bartnik and Chru\'sciel prove positive energy and parallel-spinor
rigidity with a noncompact interior
\cite[Theorem~11.2]{BartnikChrusciel2005}; Cecchini and Zeidler give a
smooth Riemannian spin formulation with arbitrary other ends and
Euclidean rigidity \cite[Theorem~B]{CecchiniZeidler2024}.
We supply the spinor argument at the rough regularity of the
three-dimensional double. In dimension four, Lesourd, Unger and Yau
supply smooth nonspin nonnegative mass for a distinguished
asymptotically Schwarzschild end, allowing arbitrary other ends
\cite[Theorem~1.2 and Definition~1.9, arXiv version~1]{LesourdUngerYau2021}.
We prove the required rigidity by deformation and volume comparison.
In the strong-decay branch, the Riemannian numerical companion
supplies the mass bound in every dimension from which the rough-double
rigidity is derived. None of these applications identifies a merely numerical
inequality with an equality theorem.

\subsection{Reading the proof}

Section~\ref{sec:setting} states the strong-decay equality theorem and
converse. Section~\ref{q:geometry:section} establishes the common
enclosure geometry and metric derivative.
Section~\ref{sec:variation} carries out the all-dimensional separation
and support localization; its smooth-dimensional shortcut is
identified after the contact argument. Section~\ref{sec:static}
contains the common stationary construction and the all-dimensional
classification, followed by the converse ADM calculation in
Section~\ref{q:converse:sec}.

For weak decay in dimension three, Sections~\ref{q3:intro:section}
and~\ref{q3:rigidity:setup} give the data and equality statements.
Variation and the static base lead to global classification and
horizon recovery in Section~\ref{q3:global:section}. A separate local
two-component spinor argument follows that main route. It examines a
family of complete conformal metrics and obtains a local obstruction
in the limit, without requiring completeness of the limiting metric.
It provides an alternative explanation of strictness for two spherical
components. The four-dimensional route begins in
Section~\ref{q4:sec:setup} and concludes with horizon recovery and the
weak-decay converse in Section~\ref{q4:rec:section}.
